% Figure: a dispatchable plant following a steep net-load ramp against its
% ramp-rate limit and minimum stable load. The net-load demand (blue) rises
% faster in the evening than the plant can follow at its rated ramp rate, so
% the plant output (orange) lags on the up-ramp and is floored at its minimum
% stable generation overnight rather than shutting down. The shaded gap is the
% unserved ramp that faster storage or another unit must cover. Values
% illustrative.
\begin{figure}[htbp]
\centering
\begin{tikzpicture}
\begin{axis}[
  width=13cm, height=8cm,
  xlabel={hour of day},
  ylabel={power (fraction of plant rating)},
  xmin=0, xmax=24, ymin=0, ymax=1.15,
  xtick={0,4,8,12,16,20,24},
  ytick={0,0.2,0.4,0.6,0.8,1.0},
  axis lines=left,
  tick label style={font=\scriptsize},
  label style={font=\small},
  legend style={font=\scriptsize, at={(0.03,0.97)}, anchor=north west, draw=none},
]
% net-load demand the plant is asked to follow: deep midday trough (solar),
% steep evening ramp, high evening peak
\addplot[engblue, very thick] coordinates {
  (0,0.55)(2,0.48)(4,0.45)(6,0.50)(8,0.60)(10,0.45)(12,0.28)
  (14,0.25)(16,0.40)(17,0.60)(18,0.85)(19,1.00)(20,0.95)
  (21,0.80)(22,0.70)(24,0.58)};
\addlegendentry{net-load demand}
% plant output: floored at minimum stable load 0.30, ramp-limited on the
% steep evening rise so it lags the demand
\addplot[engorange, very thick] coordinates {
  (0,0.55)(2,0.48)(4,0.45)(6,0.50)(8,0.60)(10,0.45)(12,0.30)
  (14,0.30)(16,0.40)(17,0.55)(18,0.72)(19,0.89)(20,0.95)
  (21,0.80)(22,0.70)(24,0.58)};
\addlegendentry{plant output (ramp- and floor-limited)}
% minimum stable load line
\addplot[enggray, thick, dashed] coordinates {(0,0.30)(24,0.30)};
\node[anchor=west, font=\scriptsize, color=enggray]
  at (axis cs:0.5,0.235) {minimum stable load};
% mark the ramp gap in the early evening
\node[anchor=west, font=\scriptsize, color=engred]
  at (axis cs:15.0,0.80) {ramp gap};
\draw[engred, -{Latex[length=2mm]}] (axis cs:16.7,0.78) -- (axis cs:18.0,0.79);
\end{axis}
\end{tikzpicture}
\caption[A dispatchable plant against its ramp and minimum-load limits]{A
dispatchable plant asked to follow the evening net-load ramp (blue) cannot
track it exactly. Two constraints bind. On the steep early-evening rise the
plant output (orange) lags the demand because its rated ramp rate limits how
fast it can raise power, and the gap between demand and output is the unserved
ramp that faster storage or a second unit must cover. Overnight and through
the midday solar trough the output is held at the minimum stable load rather
than shutting down, because cycling the unit off and on stresses it and costs
fuel and time to restart, so it runs part-loaded and inefficient through the
low hours. The steeper the net-load ramp grows with renewable penetration, the
harder both limits bind, which is the operational reason a high-renewable grid
values fast-ramping flexibility over cheap bulk energy. Quantities
illustrative.}
\label{fig:vol04ch14:ramp-cycling}
\end{figure}
